% ------------------------------------------------------------
% Page 85
% ------------------------------------------------------------

descente de Ste Enimie la plupart des sacs dégringolèrent dans le ravin. Nous ne fûmes jamais
payés.

L’usine redémarra sous un jour nouveau. Il y a lieu de dire qu’en juin les gens d’origine
italienne furent réquisitionnés par le STO. C’est ainsi que notre excellent chauffeur Persico,
bien que naturalisé français dut partir : même chose pour ce pauvre Pugnale Gino dont le fils
était infirme.

\vspace{0.45cm}

Fin juillet ou début Août après la libération d’Alès et de Nîmes ceux qui avaient pris le
maquis revinrent à l’usine. Ils avaient bien combattu dans la région du Vigan, Ganges et Saint
Hippolyte, retardant les colonnes allemandes sur la nationale 99.

\vspace{0.55cm}

\textbf{Après la libération}

\vspace{0.35cm}

Après la fin de cette période d’incertitude où tout le monde était sur le qui-vive j’interrogeais
tous ceux qui s’étaient fait inscrire sous un nom d’emprunt. Je fis un état que j’envoyais à la
Sécurité Sociale à Mende pour signaler que les cotisations versées sous le nom de M. X
concernaient en réalité M. Y ... né à ... etc... Ainsi les cotisations versées ne furent pas
perdues. Sur un effectif de 35 personnes travaillant en usine pendant cette période 1943 –
1944 la moitié était en situation irrégulière :

\vspace{0.25cm}

Evadés d’Allemagne Prisonnier de Guerre \hfill 1

-------------------S T O \hfill 1

Requis évadés ou ne s’étant pas présentés \hfill 12

Planqués avant d’être requis \hfill 6

\vspace{0.4cm}

Ces chiffres je les donne avec réserve et je les base sur des souvenirs en faisant appel à ma
mémoire. Je pense qu’ils doivent être à peu près exacts.
Travaillaient en forêt le contremaître forestier Neufeld et des équipes de bûcherons. L’effectif
forestier était d’une quinzaine. Parmi eux se trouvaient également des garçons évadés du
STO.

\vspace{0.55cm}

\textbf{L’usine de Carburants Forestiers et Meyrueis}

\vspace{0.35cm}

L’usine a rendu de très grands services à la population de Meyrueis. Rolland lorsqu’il
descendait à Millau ou à Nîmes emmenait toujours quelqu’un dans la voiture de liaison.
Lorsque les camions revenaient de livrer du charbon de bois ils aidaient au ravitaillement. Il
s’était créé également des liens de solidarité parmi tout le personnel. Il n’y eut aucune
dénonciation bien qu’il y eut dans l’usine, 3 ouvriers qui, par bêtise, s’étaient fait inscrire à la
Milice, soi-disant pour ne pas partir en Allemagne bien qu’à l’usine existât cette garantie !!! ;
mais n’en bénéficièrent pas les étrangers !

\vspace{0.55cm}

\textbf{La page est tournée.}

\vspace{0.35cm}

Notre usine des Carburants Forestiers a tenu encore pendant 5 ans après la libération. Elle
s’était orientée vers d’autres fabrications mais par suite de l’éloignement et de complications
administratives elle fut obligée de déposer son bilan.
